% Mục Kết quả, mutex lock (RQ2, RQ3). Nguồn: results/apple-m5/lock.csv, lock_yield.csv,
% sleep_latency.csv (trung vị 7 lần lặp); bảng locktime, locktime-ow*, locksingle, sleeplatency,
% lockunit trong report/figures/apple-m5/tables.tex. Số không có trong bảng in ra (độ chậm, độ công
% bằng, ns mỗi critical section, tỉ số backoff) tính bằng scripts/report_numbers/rlock_analyze.py và
% scripts/report_numbers/rlock_analyze2.py; hai script nạp scripts/plot.py và dùng đúng aggregate/unit_time/bound của nó, nên trùng
% với số vẽ trên hình. Số ghi "lock.csv: ..." tính thẳng từ CSV: ns mỗi critical section = trung vị
% seconds/ops*1e9, độ công bằng = trung vị thr_min_s/thr_max_s của cột extra. Nguồn từng số ghi trong comment ngay cạnh.

Hai đại lượng dùng trong mục này là độ chậm và độ công bằng, định nghĩa ở mục Phương pháp đo.

% ================================================================================================
\subsubsection{RQ2: Lock mở rộng thế nào theo số thread, tranh chấp và độ dài critical section?}

\begin{rqanswer}\textbf{Trả lời RQ2.} TTAS mở rộng tốt nhất trong ba lock tự viết. Ở 20 thread
($N = 1009$, $\mathrm{OW} = 0$), nó chậm 2,48 lần cận dưới. MCS gãy mạnh tại số core (13 lần ở 10
thread, 7 lần khi tải nền thấp) và sụp khi vượt số core (1\,700--6\,900 lần). Các ô đó đo với số vòng lặp giảm 100 lần và chỉ cho bậc độ
lớn. Lock thư viện không sụp. \PCname{std_mutex} và \PCname{pthread_mutex} chậm hơn spinlock khi
critical section dài cỡ 0,5~$\mu$s (3,9--4,0 lần cận dưới ở 20 thread, so với 2,48 của TTAS). Riêng
\PCname{os_unfair_lock} giữ 1,6--1,7 lần. Khi critical section chỉ vài nanosecond, cả ba lock thư viện
nhanh hơn spinlock không backoff từ 3 thread. Số liệu của \PCname{apple-m5}. RQ6 kiểm lại trên hai máy còn lại.\end{rqanswer}
% nguồn: scripts/report_numbers/rlock_analyze.py, độ chậm N=1009 OW=0: p=20 ttas 2,48, std_mutex 3,93, pthread_mutex 3,97,
% os_unfair_lock (p=2..20) 1,59..1,69; mcs p=10 13,0, p=12/16/20 1,70e3/5,16e3/6,88e3 (đo giảm vòng).
% N=11 OW=0 p>=3: std_mutex 4,08..5,81, pthread_mutex 3,68..5,14, os_unfair_lock 2,87..5,02; tas 5,18..80,5, ttas 7,28..28,2,
% mcs 18,7..101 (p<=10). Ở p=2 os_unfair_lock 5,57 > tas 4,39, nên câu chỉ đúng từ 3 thread.

\PCfigure{apple-m5}{lock_slowdown}

\emph{Bằng chứng.} Ở $\mathrm{OW} = 0$, ba spinlock tự viết chậm không quá 1,24 lần cận dưới tới 4
thread (Hình~\ref{fig:apple-m5:lock_slowdown}). Chúng gãy từ 4 lên 6 thread (TAS 1,18 lên 1,50, TTAS
1,08 lên 1,31, MCS 1,24 lên 2,98), tức tại số P-core chứ không tại số core. MCS gãy mạnh nhất. Từ 6
thread nó chậm nhất trong ba spinlock tự viết, và từ 8 thread chậm nhất trong mọi lock (4,31 so với 3,72
của \PCname{std_mutex}). Nó chậm 13 lần ở đúng 10 thread rồi sụp hẳn khi vượt số core
(Bảng~\ref{tab:apple-m5:locktime}). TAS và TTAS chỉ chậm dần, tới 3,48 và 2,48 lần ở 20 thread. Lock
thư viện đã chậm 2,5--2,9 lần từ 3--4 thread (\PCname{os_unfair_lock} 1,6 lần từ 2 thread), rồi tăng
chậm và gần như phẳng từ 8 thread (3,5--4,0 lần).
% nguồn: scripts/report_numbers/rlock_analyze.py, độ chậm N=1009 OW=0 p=4/6: tas 1,18/1,50, ttas 1,08/1,31, mcs 1,24/2,98;
% mcs p=8 4,31, std_mutex p=8 3,72; mcs p=10 13,0; p=20 tas 3,48, ttas 2,48.
% std_mutex p=3/4 2,79/2,93; pthread_mutex 2,45/2,86; p=8..20 std_mutex 3,72..3,93, pthread_mutex 3,54..3,97;
% os_unfair_lock (p=2..20) 1,59..1,69.

\PCtable{apple-m5}{locktime}

Ô $\mathrm{OW} = 15$ là một thí nghiệm đối chứng, vì với $p \le 15$ lock bận tối đa $p/16$ thời gian.
Vậy mà ở 6 thread cả mười hai lock, kể cả lock thư viện, cùng chậm 1,53--2,03 lần cận dưới.
% nguồn: scripts/report_numbers/rlock_analyze.py, độ chậm N=1009 OW=15 p=6: 1,53 (mcs) .. 2,03 (os_unfair_lock).

Từ $p = 1+\mathrm{OW}$, lock phải bận liên tục thì mới đạt cận dưới, nên mỗi lần chuyển giao chậm
đều cộng thẳng vào thời gian chạy. Lock thư viện cho thấy điểm này rõ nhất. Ở $\mathrm{OW} = 1$ và 2
thread, \PCname{std_mutex} và \PCname{pthread_mutex} đã chậm 2,9 lần cận dưới (\PCname{os_unfair_lock}
1,7 lần), còn ba spinlock tự viết chậm 1,11--1,27 lần. Ở $\mathrm{OW} = 3$, ba lock thư viện chậm
1,34--1,83 lần ở 3 thread và 2,64--2,86 lần ở 4 thread, còn spinlock 1,11--1,13 và 1,31--1,37 lần. Nhóm cho rằng
thread chờ của lock thư viện ngủ, nên mỗi lần chuyển giao trả thêm một lần đánh thức qua kernel (ước
lượng ở dưới), và lock thư viện bão hoà sớm hơn spinlock. Với $\mathrm{OW} = 7$ và 15, chỗ gãy từ 4 lên
6 thread tới trước điểm bão hoà (TAS 1,09 lên 1,92 và 1,05 lên 1,59), và mọi lock tăng cùng nhau.
% nguồn: scripts/report_numbers/rlock_analyze.py, độ chậm N=1009 OW=1 p=2: std_mutex 2,86, pthread_mutex 2,91,
% os_unfair_lock 1,71; tas/ttas/mcs 1,12/1,11/1,27. OW=3 p=3: std_mutex 1,83, pthread_mutex 1,68, os_unfair_lock 1,34,
% tas/ttas/mcs 1,13/1,11/1,12; p=4: std_mutex 2,64, pthread_mutex 2,86, os_unfair_lock 2,84, tas/ttas/mcs 1,31/1,37/1,37.
% tas OW=7 p=4/6 1,09/1,92; OW=15 p=4/6 1,05/1,59.

Khi vượt số core ở $\mathrm{OW} = 15$, riêng phần việc đã vượt sức máy. Tổng $p(1+\mathrm{OW})\,I$ lời
gọi \texttt{is\_prime} chia cho 10 core cần ít nhất $p(1+\mathrm{OW})\,I\,t/10$, tức 1,6 lần cận dưới ở
16 và 20 thread, trước khi tính chi phí lock. TAS, TTAS và lock thư viện cùng nằm trong 4,9--6,0 lần ở
16--20 thread. Nhóm cho rằng phần còn lại đến từ E-core chậm hơn và từ việc chia lượt lập lịch.
% nguồn: số học từ định nghĩa cận dưới max(p,1+OW) I t (mục Phương pháp đo): p(1+OW)/(10 max(p,1+OW)) = 16*16/160 = 1,6
% (p=16), 20*16/200 = 1,6 (p=20). scripts/report_numbers/rlock_analyze.py, độ chậm N=1009 OW=15 p=16/20: tas 5,28/6,02,
% ttas 4,93/5,82, std_mutex 4,92/4,95, pthread_mutex 4,88/5,07, os_unfair_lock 5,07/5,06.

Với $N = 10\,007$ (khoảng 5,2~$\mu$s) TAS và TTAS gần như trùng nhau
(Hình~\ref{fig:apple-m5:lock_cs_length}, phụ lục). Với $N = 11$ (4,6~ns), ở 20 thread TAS chậm 80,5 lần cận
dưới, TTAS 28,2 lần, còn lock thư viện chỉ khoảng 5 lần.
% nguồn: scripts/report_numbers/rlock_analyze.py, độ chậm OW=0: N=10007 p=20 tas 2,55, ttas 2,53; N=11 p=20 tas 80,5, ttas 28,2,
% std_mutex 5,71, pthread_mutex 4,97, os_unfair_lock 5,02.

\emph{Vì sao.} Nhóm cho rằng chỗ gãy ở 4 thread đến từ máy, không từ thuật toán. Ở ô đối chứng, mọi
lock cùng chậm như nhau. Độ công bằng ở 6 thread là 0,84--0,98, tức mọi thread cùng chậm đi. Điều này
phù hợp với việc scheduler luân chuyển thread giữa P-core và E-core chậm hơn, hoặc với tải nền (xem mục
Các yếu tố đe dọa độ tin cậy). Trên \PCname{C4M7Y670NT} (6 P-core, 12 core), điểm gãy dời sang 6 lên 8
thread và MCS sụp khi vượt 12 thread (mục So sánh giữa các máy).
% nguồn: scripts/report_numbers/rlock_analyze.py, fairness N=1009 OW=15 p=6: 0,84 (mcs) .. 0,98.
% nguồn: results/C4M7Y670NT/lock.csv (độ chậm tính như rlock_analyze.py: trung vị / (iters t p), t từ lần chạy 1 thread
% nhanh nhất), N=1009 OW=0: tas p=6/8 1,15/1,53; mcs p=6/8 1,23/2,58 (bảng locktime: 1,08/0,937, 1,15/0,937), p=12 4,99, p=16 2\,070.

MCS gãy mạnh ngay tại số core và sụp khi vượt số core, sớm hơn mức mà riêng lưu lượng cache coherence
giải thích được. Nguyên nhân nằm ở cách trao lock. MCS trao lock cho đúng một thread kế tiếp theo thứ tự FIFO (vào trước được phục
vụ trước). Nếu thread đó không có core, lock nằm yên cho tới khi scheduler cho nó chạy lại, và cả hàng
chờ mất một lượt lập lịch. TAS và TTAS không bị như vậy, vì thread nào đang chạy cũng lấy được lock
rảnh. Độ công bằng của MCS bằng 1 tới 10 thread (Hình~\ref{fig:apple-m5:lock_fairness}), tức lock luôn
đi đúng thứ tự hàng đợi.
% nguồn: scripts/report_numbers/rlock_analyze.py, fairness mcs N=1009 OW=0 p=2..10: 1,00.

Chương trình \texttt{handoff\_gap} kiểm cơ chế này. Nó chạy đúng vòng lặp của phần lock ($N = 1009$,
$\mathrm{OW} = 0$) cho MCS, TAS, TTAS và \PCname{std_mutex} ở 4, 8, 10 và 12 thread. Mỗi lần chạy có
khoảng 40\,000 critical section (3\,000 cho MCS ở 12 thread), và mỗi cấu hình chạy 5 lần theo thứ tự
xáo trộn. Thread nhả lock ghi lại thời điểm nhả, và thread lấy lock kế tiếp đo \emph{khoảng trống trao
lock}, tức thời gian từ lúc lock được nhả tới lúc nó có chủ mới. Ở 4 thread, trung vị của khoảng trống
dài nhất trong một lần chạy là 6--22~$\mu$s với cả bốn lock. Vì vậy một khoảng trống dài hơn
100~$\mu$s được coi là lock đã chờ một thread không có core. Ở 4 thread, chỉ một trong 20 lần chạy có
khoảng trống như vậy (MCS, hai lần trao lock).
% nguồn: experiments/run_extra.py (part_a: total = 3000 nếu mcs và p>=12, ngược lại 40000; iters = total // p;
% 5 rep, random.Random(SEED+rep).shuffle). handoff.csv: handovers 39\,999 (p=4,8,10), 39\,995 (p=12), 2\,999 (mcs p=12).
% max_gap_us p=4, trung vị 5 lần (extra_analyze.py): ttas 6, tas 10, std_mutex 14, mcs 22; lần lớn nhất: mcs 199,5
% (rep 3, long_handovers 2), tas 79,1, ttas 69,0, std_mutex 18,1.

Bảng~\ref{tab:handoff} cho thấy MCS ở 10 thread chỉ có 0,25\,\% lần trao lock dài như vậy, mỗi lần
trung bình 0,79~ms, nhưng chúng chiếm 47\,\% thời gian chạy. Ở 12 thread, 27\,\% lần trao lock là lần
dài và chúng chiếm 99\,\% thời gian. TAS, TTAS và \PCname{std_mutex} gần như không có lần trao lock dài
nào. Số lần chuyển ngữ cảnh không tự nguyện của tiến trình MCS tăng từ 60 ở 4 thread lên 1\,970 ở 10
thread và 23\,044 ở 12 thread. Phép đo này chạy khi load average khoảng 2,8, nên ngay ở 10 thread đã có
thread bị tước core. Đối chứng là TAS ở 10 thread. Tiến trình TAS cũng bị chuyển ngữ cảnh 1\,022 lần,
cùng bậc với MCS, nhưng trung vị số lần trao lock dài bằng 0, vì thread nào đang chạy cũng lấy được lock
rảnh. Phép đo này không kiểm điểm gãy từ 4 lên 6 thread, nên cơ chế của điểm gãy đó vẫn là giả thuyết.
% nguồn: scripts/report_numbers/extra_analyze.py (results/apple-m5/extra/handoff.csv, trung vị 5 lần chạy):
% mcs p=10 long 0,245 %, mean_long 792 us, time_in_long 47,3 %; p=12 26,8 % / 2\,233 us / 99,4 %;
% nivcsw mcs p=4/10/12: 60 / 1\,970 / 23\,044; load1 2,8--2,9.
% tas p=10: nivcsw 1022/2317/466/753/1024 (trung vị 1\,022), long_handovers 0/8/2/0/0 (trung vị 0).

\begin{table}[H]\centering\small
\begin{tabular}{@{}lcccc@{}}
\toprule
 & 4 thread & 8 thread & 10 thread & 12 thread \\
\midrule
\PCname{mcs}       & 0 / 0    & 0,057 / 8,3 & 0,25 / 47 & 27 / 99 \\
\PCname{tas}       & 0 / 0    & 0 / 0       & 0 / 0     & 0,003 / 0,3 \\
\PCname{ttas}      & 0 / 0    & 0 / 0       & 0 / 0     & 0 / 0 \\
\PCname{std_mutex} & 0 / 0    & 0 / 0       & 0 / 0     & 0 / 0 \\
\bottomrule
\end{tabular}
\caption{Lần trao lock dài hơn 100~$\mu$s trên \PCname{apple-m5}, $N = 1009$, $\mathrm{OW} = 0$: tỉ lệ
trên mọi lần trao lock (\%) / phần thời gian chạy nằm trong các lần đó (\%). Trung vị của 5 lần chạy.}
\label{tab:handoff}
\end{table}
% nguồn: scripts/report_numbers/extra_analyze.py (cột long% và time_in_long%).

Cùng giả thuyết phù hợp với việc MCS gãy mạnh hơn TAS và TTAS từ 4 lên 6 thread. MCS buộc lock đi qua
mọi thread, kể cả thread đang chạy trên E-core hay đang chờ core. TAS và TTAS cho phép thread vừa nhả
lấy lại lock. Ở 4 thread độ công bằng của chúng là 0,49 và 0,26, của MCS là 1,00.
% nguồn: scripts/report_numbers/rlock_analyze.py, fairness N=1009 OW=0 p=4: tas 0,49, ttas 0,26, mcs 1,00.

Lock thư viện cho thread chờ ngủ trong kernel, nên thread chờ không giành core của chủ lock và lock
không sụp. Thread đang ngủ cũng không đọc cache line của lock, nên thêm thread chờ gần như không thêm
chi phí. Một critical section của \PCname{std_mutex} tốn 1\,901~ns ở 8 thread và 2\,007~ns ở 20 thread.
% nguồn: scripts/report_numbers/rlock_analyze.py, ns/CS std_mutex N=1009 OW=0 p=8/20: 1\,901 / 2\,007.

Nhóm giả thuyết rằng mỗi lần đánh thức một thread qua kernel tốn một chi phí cố định cỡ 1,5~$\mu$s. Con
số này là ước lượng suy ra từ chính số đo lock. Báo cáo không có phép đo đánh thức riêng. Ở
$N = 1009$, mỗi critical section của \PCname{std_mutex} ở 20 thread tốn thêm 1\,459~ns so với khi không
tranh chấp (548~ns ở 1 thread, Bảng~\ref{tab:apple-m5:locksingle}, phụ lục). Ở $N = 10\,007$, mức tăng là 1\,469~ns, gần như không đổi dù phần việc dài gấp 10 lần. Mức
tăng đó bằng khoảng 0,28 lần một lời gọi \texttt{is\_prime(10\,007)}, khớp với mức 1,25--1,41 lần cận
dưới của lock thư viện từ 4 thread ở $N = 10\,007$. Nhóm cho rằng \PCname{pthread_mutex} (và \PCname{std_mutex} dựng
trên nó) trên macOS trao lock cho thread được đánh thức, nên gần như mỗi lần chuyển giao trả trọn một
lần đánh thức. Tới 10 thread, độ công bằng 0,88--0,97 của hai lock này phù hợp với điều đó. Khi vượt
số core, nó giảm dần xuống 0,78--0,79 ở 20 thread. Nhóm cho rằng khi đó scheduler chia core không đều
giữa các thread, độc lập với lock. Ở 2 thread mức tăng mới là 316~ns. Có thể thread chờ duy nhất nhiều lần chưa kịp ngủ
khi lock được nhả.
% nguồn: scripts/report_numbers/rlock_analyze.py, ns/CS std_mutex N=1009 OW=0 p=20: 2\,007 ns; locksingle 548; 2\,007 - 548 = 1\,459;
% p=2: 864,3 - 548,1 = 316. lock.csv (N=10007 OW=0, ns/CS): std_mutex p=1 5\,377, p=20 6\,847; 6\,847 - 5\,377 = 1\,469;
% 1\,469 / 5\,163 = 0,28. rlock_analyze.py độ chậm N=10007 OW=0 p=4..20: std_mutex 1,25..1,33, pthread_mutex 1,28..1,41 (1,41 ở p=16),
% os_unfair_lock 1,25..1,30. fairness N=1009 OW=0: std_mutex p=2..10 0,92..0,97, p=12/16/20 0,89/0,83/0,78;
% pthread_mutex p=2..10 0,88..0,96, p=12/16/20 0,87/0,86/0,79.

Tài liệu của \PCname{os_unfair_lock} (\texttt{os/lock.h}) ghi rằng thread vừa nhả có thể lấy lại lock
trước khi thread được đánh thức kịp thử. Phần lớn critical section vì thế không phải trả lần đánh thức
nào. Một critical section chỉ tốn thêm 302--351~ns từ 2 tới 20 thread, và độ công bằng giảm từ 0,97 ở 3
thread xuống 0,20 ở 20 thread.
% nguồn: MacOSX.sdk/usr/include/os/lock.h, đoạn @discussion của os_unfair_lock. scripts/report_numbers/rlock_analyze.py,
% ns/CS os_unfair_lock N=1009 OW=0: p=1 510,6; p=2..20 812,3..861,6 (nhỏ nhất p=2, lớn nhất p=16), trừ 510,6 = 301,7..351,0;
% fairness p=3 0,97, p=20 0,20.

Ở $N = 11$, \PCname{std_mutex} tốn thêm 19,7~ns mỗi critical section. Nếu mỗi lần đánh thức tốn
1\,459~ns như ở $N = 1009$, một lần đánh thức được chia cho khoảng 74 critical section. Nhóm cho rằng trong lúc thread được đánh thức chưa kịp chạy, thread vừa
nhả lock lấy lại lock và chạy tiếp các critical section của mình với cache line còn nóng. Chuỗi này kéo
dài khoảng 1,9~$\mu$s, cùng bậc với lần ngủ ngắn nhất của \texttt{sleep\_for} khi yêu cầu 100~ns
(2\,540~ns, Bảng~\ref{tab:apple-m5:sleeplatency}, phụ lục). Chi phí thêm mỗi critical section nhỏ hơn
một critical section của TAS khi tranh chấp (39,5~ns ở 4 thread, gồm 4,6~ns của
\texttt{is\_prime(11)}), nên phần lớn critical section không thể có chuyển giao. Độ công bằng 0,92 của
\PCname{std_mutex} ở 20 thread không chống lại giả thuyết. Thước đo này không thấy các đợt độc chiếm
luân phiên (mục Các yếu tố đe dọa độ tin cậy).
% nguồn: scripts/report_numbers/rlock_analyze.py, ns/CS std_mutex N=11 OW=0 p=20: 26,24; locksingle N=11 std_mutex 6,50;
% 26,24 - 6,50 = 19,74; 1\,459 / 19,74 = 74; 74 * 26,24 ns = 1,94 us. sleeplatency (results/apple-m5/sleep_latency.csv):
% sleep_for yêu cầu 100 ns, nhỏ nhất 2\,541 ns, trung vị 2\,667 ns. tas N=11 p=4: 39,53 ns/CS.
% fairness std_mutex N=11 OW=0 p=20: 0,92 (scripts/report_numbers/rlock_analyze2.py).

TAS thua TTAS nhiều nhất khi lock đổi chủ ở mỗi lượt. Phép nhả lock của TAS phải tranh cache line với
các \texttt{exchange} liên tục của thread chờ, còn thread TTAS chờ bằng phép đọc. Ở $\mathrm{OW} = 1$ và
8 thread, TAS đạt 0,467 triệu lần lấy lock mỗi giây, TTAS 1,41 triệu. Khi không có việc ngoài, lock chủ
yếu ở lại với thread vừa nhả. Độ công bằng của TAS 10 thread là 0,32 ở $\mathrm{OW} = 0$ và 0,91 ở
$\mathrm{OW} = 3$. Vì vậy TAS ở $\mathrm{OW} = 1$ và 3 chậm hơn ở $\mathrm{OW} = 0$ (6,79 và 7,17 so với
3,48 lần cận dưới ở 20 thread).
% nguồn: scripts/report_numbers/rlock_analyze.py, fairness N=1009 tas p=10: OW=0 0,32; OW=3 0,91. Thông lượng OW=1 p=8
% ttas 1,41 M/s, tas 0,467 M/s (locktime-ow1: 0,975 s và 2,94 s). Độ chậm tas p=20: OW=0 3,48, OW=1 6,79, OW=3 7,17.

Độ dài critical section quyết định tỉ lệ giữa chi phí chuyển giao và phần việc
(Hình~\ref{fig:apple-m5:lock_cs_length}, phụ lục). Ở $N = 11$ và 4 thread, một critical section của TAS tốn
39,5~ns, gấp khoảng 8,6 lần lời gọi \texttt{is\_prime(11)}. MCS chậm nhất ở điểm này (86~ns mỗi
critical section, 18,7 lần cận dưới). Nhóm cho rằng thứ tự FIFO buộc mọi critical section đổi chủ, còn
TAS và TTAS cho phép thread vừa nhả lấy lại lock thành từng đợt. Ở $N = 10\,007$, chi phí chuyển giao của
TAS chưa tới 1\,\% của 5,2~$\mu$s. Tới 4 thread mọi lock nằm trong 1,0--1,3 lần cận dưới. Từ 6 thread
TAS và TTAS tách dần (1,56 và 1,57 lần ở 10 thread). Khi vượt số core, khác biệt quay lại rõ. TAS và
TTAS chậm 2,5 lần ở 20 thread, lock thư viện 1,3 lần, còn MCS sụp. Nhóm cho rằng thread chờ quay giành core của
chủ lock, còn thread chờ của lock thư viện thì ngủ.
% nguồn: scripts/report_numbers/rlock_analyze.py ns/CS tas N=11 OW=0 p=4: 39,53; t(N=11) 4,60 ns (lockunit); 39,53/4,60 = 8,6;
% 39,53/5\,163 = 0,77 %. ns/CS mcs N=11 OW=0 p=4: 86,06; độ chậm 18,7. Độ chậm N=10007 OW=0 p=1..4: 1,00..1,32;
% p=10: tas 1,56, ttas 1,57; p=20: tas 2,55, ttas 2,53, std_mutex 1,33, pthread_mutex 1,28, os_unfair_lock 1,30, mcs 642.

Khi vượt số core, một critical section của MCS tốn gần như cùng thời gian ở $N = 1009$ và
$N = 10\,007$ (3,5 và 3,3~ms ở 20 thread), dù phần việc dài gấp 10 lần. Chi phí đó nằm ngoài phần việc.
Độ chậm của MCS ở $N = 10\,007$ (642 lần) thấp hơn ở $N = 1009$ (6\,880 lần) chỉ vì cận dưới lớn hơn.
% nguồn: lock.csv (mcs OW=0 p=20, ns/CS): N=1009 3,51e6, N=10007 3,31e6. rlock_analyze.py độ chậm mcs p=20 OW=0:
% N=1009 6\,880, N=10007 642.

\emph{Độ vững.} Cả bảy lần chạy MCS 10 thread đều chậm 7,1--15,5 lần cận dưới. Khi đo lại với load
average khoảng 3 thay vì khoảng 25, chỉ MCS ở 10 thread nhanh lên rõ, gần gấp đôi. Thời gian của các
lock khác bằng 0,83--1,05 lần thời gian ở đợt đo chính (mục Các yếu tố đe dọa độ tin cậy). Điều này phù hợp với cơ chế trên.
% nguồn: scripts/report_numbers/extra_analyze.py (lowload.csv): mcs p=10 low/main 0,53 (3\,545 / 6\,635 ns/CS);
% các dòng còn lại (kể cả mcs p<=8) 0,83 (tas 10t) .. 1,05 (std_mutex 4t).

Các điểm MCS vượt số core ở $N = 11$ không đo được sự sụp đổ, vì sau lần vượt thời hạn số vòng lặp giảm
tới mức hàng đợi không kịp hình thành. Ở $\mathrm{OW} \ge 1$, một số ô MCS vượt số core cho
kết quả hai cực, nên chỉ khoảng min--max đáng tin. Ở $\mathrm{OW} \ge 7$, trung vị ở 20 thread rơi vào cực
nhanh. Ở 16 thread, $\mathrm{OW} = 15$, bốn lần chạy xong dưới 10~ms, ba lần mất 2,0--10,6~s. Nhóm cho
rằng khi phần lớn thời gian nằm ngoài lock, hàng chờ thường ngắn. Lần chạy nào không có lượt trao lock
cho thread mất core thì xong nhanh, lần chạy nào có thì kéo dài nhiều giây.
% nguồn: lock.csv, mcs p=10 OW=0 N=1009 iters=34\,366: 1,241..2,727 s, cận dưới 0,1755 s.
% 6 dòng timeout: mcs N=11 OW=0/3 p=12/16/20, iters 4\,522/3\,391/2\,713; lần sau 45/33/27 vòng, ~0,1 ms.
% Hai cực (rlock_analyze.py, mcs N=1009, min..max s): OW=1 p=16 0,0071..7,96; OW=3 p=12 0,024..3,86, p=16 0,0057..8,75,
% p=20 0,0055..18,3; OW=7 p=20 trung vị 0,020 (0,0052..13,5); OW=15 p=20 trung vị 0,0075 (0,0066..17,4).
% mcs p=16 OW=15: 4 lần < 10 ms, 3 lần 2,0..10,6 s.

% ================================================================================================
\subsubsection{RQ3: Backoff giúp hay hại, và lock nhanh nhất có công bằng không?}

\begin{rqanswer}\textbf{Trả lời RQ3.} Backoff giúp TAS và gần như không đổi gì với TTAS. Ở 20 thread
($N = 1009$, $\mathrm{OW} = 0$), backoff bằng \texttt{sleep\_for} đưa TAS từ 12,2~s xuống 3,78~s
(thời gian quy về cùng số vòng lặp). Lợi ích đến từ việc lock gần như không đổi chủ, nên lock nhanh
nhất lại kém công bằng nhất. Khi có việc ngoài critical section và tranh chấp vừa phải, backoff làm
chậm. Hại lớn nhất là ở \PCname{ttas_bo_sleep} khi vượt số core với $\mathrm{OW} = 3$, chậm hơn TTAS
6,1--6,8 lần. Số liệu của \PCname{apple-m5}. RQ6 kiểm lại trên hai máy còn lại.\end{rqanswer}
% nguồn: locktime p=20: tas 12,2; tas_bo_sleep 3,78. rlock_analyze.py tỉ số ttas_bo_sleep/ttas N=1009 OW=3
% p=16/20: 6,13/6,78.

\PCfigure{apple-m5}{lock_backoff}

\emph{Bằng chứng.} Ở $\mathrm{OW} = 0$ (Hình~\ref{fig:apple-m5:lock_backoff}, cột trái),
\PCname{tas_bo_sleep} và \PCname{tas_bo_yield} chỉ mất 0,31--0,32 lần thời gian của TAS ở 20 thread,
còn \PCname{tas_bo_spin} mất 0,75 lần. Các biến thể của TTAS nằm trong 0,87--1,15 lần TTAS. Ở
$\mathrm{OW} = 3$ và 4--6 thread, backoff làm TAS chậm hơn (\PCname{tas_bo_sleep} 1,32--1,33 lần,
\PCname{tas_bo_spin} 1,29--1,49 lần). Ở 20 thread, \PCname{tas_bo_sleep} và \PCname{tas_bo_yield} lại
nhanh hơn (0,69 và 0,55 lần TAS), còn \PCname{tas_bo_spin} vẫn chậm hơn (1,18 lần). Trên
\PCname{LAP1184}, cũng ở $\mathrm{OW} = 3$, sleep backoff làm TAS chậm hơn cả ở số thread lớn (RQ6). Ở
$\mathrm{OW} = 15$, mọi biến thể nằm trong 0,90--1,26 lần TAS tới 12 thread. Ở 20 thread
\PCname{tas_bo_sleep} và \PCname{tas_bo_yield} nhanh hơn (0,73 và 0,63 lần).
% nguồn: locktime p=20: 3,78/12,2 = 0,31; 3,95/12,2 = 0,32. scripts/report_numbers/rlock_analyze.py tỉ số so với tas
% N=1009: OW=0 p=20 tas_bo_spin 0,745; tỉ số ttas_bo_*/ttas OW=0: 0,87..1,15. OW=3: tas_bo_sleep p=4/6 1,32/1,33,
% tas_bo_spin p=4/6 1,29/1,49; p=20 tas_bo_sleep 0,688, tas_bo_yield 0,549, tas_bo_spin 1,18.
% OW=15: p<=12 0,90..1,26; p=20 tas_bo_sleep 0,729, tas_bo_yield 0,633.
% LAP1184 (review: độ chậm N=1009 OW=3, rlock_analyze cho LAP1184): tas_bo_sleep/tas p=3..16 1,53..2,04, p=20 1,40, p=24 1,11.
% Ở OW=0 trên LAP1184 tas_bo_sleep vẫn nhanh hơn tas ở số thread lớn (1,21 so với 2,24 ở 24 thread), nên câu chỉ nói OW=3.

\PCfigure{apple-m5}{lock_fairness}

Ở $\mathrm{OW} = 0$, MCS giữ độ công bằng 1 tới 10 thread (Hình~\ref{fig:apple-m5:lock_fairness}). Độ
công bằng của \PCname{tas_bo_sleep} bằng đúng $1/p$, tới 0,05 ở 20 thread, tức là các thread chạy xong
lần lượt từng thread một. TAS và TTAS cũng thiên vị thread vừa nhả lock. TTAS thiên vị hơn TAS (0,26 so
với 0,49 ở 4 thread, 0,05 so với 0,11 ở 20 thread). Nhóm cho rằng phép nhả của TTAS không phải tranh
cache line với \texttt{exchange} nào, nên thread vừa nhả lấy lại lock trong lúc các thread chờ còn đang
đọc lại line. \PCname{os_unfair_lock}, lock thư viện nhanh nhất, kém công bằng nhất trong ba lock thư
viện (0,20 so với 0,78--0,79 ở 20 thread), vì lý do đã nêu ở RQ2.
% nguồn: scripts/report_numbers/rlock_analyze.py, fairness N=1009 OW=0: mcs p=2..10 1,00; tas_bo_sleep p=2/4/10/20 0,50/0,25/0,10/0,05;
% tas p=4/20 0,49/0,11; ttas p=4/20 0,26/0,05; os_unfair_lock p=20 0,20; std_mutex 0,78, pthread_mutex 0,79.

Ở $\mathrm{OW} = 3$, mọi lock gần như công bằng tới 4 thread. Từ 6 thread, $p$ vượt
$1+\mathrm{OW}$ nên luôn có thread chờ, và cùng sự thiên vị đó làm độ công bằng của TTAS giảm từ 0,98 ở 4 thread xuống
0,87 ở 6 thread và 0,55 ở 20 thread.
% nguồn: scripts/report_numbers/rlock_analyze.py, fairness N=1009 OW=3 ttas p=4/6/20: 0,98/0,87/0,55.

Khi $\mathrm{OW} > 0$, thread vừa nhả đi gọi \texttt{is\_prime} thay vì lấy lại lock ngay. Vì vậy
\PCname{os_unfair_lock} gần công bằng ở $\mathrm{OW} = 3$ và 15 (0,87--1,00), dù ở $\mathrm{OW} = 0$
nó chỉ còn 0,20 ở 20 thread.
% nguồn: lock.csv (N=1009, trung vị thr_min_s/thr_max_s), os_unfair_lock p=2..20: OW=3 0,87..1,00; OW=15 0,96..0,99;
% OW=7 0,94..1,00. Ngoại lệ OW=1: 0,99 ở 2 thread giảm còn 0,57 ở 20 thread (OW=1 không có trên hình).

\emph{Vì sao.} Các biến thể \PCname{tas_bo_*} lùi lại sau mọi lần \texttt{exchange} thất bại, kể cả khi
lock chỉ đang bị giữ, nên thread chờ gần như luôn ngủ hoặc nhường core. Các biến thể \PCname{ttas_bo_*}
chỉ lùi lại sau một lần \texttt{exchange} thua và vẫn đọc liên tục như TTAS (mục Cài đặt). Một lần
\texttt{sleep\_for} yêu cầu 100~ns thật ra ngủ trung vị 2\,670~ns. Ở $\mathrm{OW} = 0$, chủ lock nhả
rồi lấy lại lock ngay ở lượt kế tiếp của vòng lặp. Thread vừa thức dậy gần như luôn thấy lock bận và
ngủ tiếp với thời gian yêu cầu gấp đôi, tới khoảng 100~$\mu$s (\texttt{MAX\_DELAY} trong
\texttt{backoff.hpp}). Vì vậy một thread giữ lock tới hết vòng lặp của nó, rồi mới tới thread khác. Độ
công bằng đúng bằng $1/p$ phù hợp với điều này.
% nguồn: sleeplatency: sleep_for 100 ns trung vị 2\,667 ns. include/backoff.hpp: SleepBackoff bắt đầu
% PC_SLEEP_MIN_NS = 100, nhân đôi khi delay < PC_SLEEP_MAX_NS = 100000 ("MAX_DELAY of the assignment text").
% fairness tas_bo_sleep N=1009 OW=0 p=2..20: 0,50/0,34/0,25/0,16/0,12/0,10/0,08/0,06/0,05 (rlock_analyze.py).

\PCname{tas_bo_spin} chỉ giãn các lần thử. Thread chờ vẫn giữ core và vẫn gọi \texttt{exchange} sau
mỗi cửa sổ quay (tối đa 1\,024 lần \texttt{cpu\_relax()}, \texttt{backoff.hpp}). Nhóm cho rằng vì vậy
lợi ích của nó nhỏ hơn, và khi vượt số core nó còn chiếm core mà chủ lock cần.
% nguồn: include/backoff.hpp, SpinBackoff: cửa sổ ngẫu nhiên, nhân đôi tới PC_SPIN_MAX = 1024.

Khi $\mathrm{OW} > 0$, lock rảnh trong lúc thread chờ còn ngủ ít nhất 2,5~$\mu$s. Với
\PCname{tas_bo_spin}, cửa sổ quay lớn nhất kéo dài khoảng 10~$\mu$s (1\,024 lần \texttt{cpu\_relax()},
mỗi lần trung vị 9,6~ns), và lock cũng rảnh trong lúc đó. Ở tranh chấp vừa phải, đó là thời gian mất.
Ở số thread lớn, luôn có thread sẵn sàng lấy lock. Ở $\mathrm{OW} = 3$ và 15, khi vượt số core, thread
ngủ hoặc nhường core để lại core cho thread đang làm việc ngoài critical section. Nhóm cho rằng đó là lý
do \PCname{tas_bo_sleep} và \PCname{tas_bo_yield} nhanh hơn TAS ở 20 thread, còn \PCname{tas_bo_spin},
vẫn giữ core, thì không.
% nguồn: sleeplatency: sleep_for 100 ns nhỏ nhất 2\,541 ns; cpu_relax trung vị 9,6 ns; 1\,024 * 9,6 = 9,8 us.
% Tỉ số so với tas p=20: OW=3 tas_bo_sleep 0,688, tas_bo_yield 0,549, tas_bo_spin 1,18; OW=15 0,729 / 0,633 / 1,11.

MCS chia lượt đều khi mọi thread được chạy, và cả khi vượt số core nếu $\mathrm{OW} \le 3$ (độ công bằng
0,84--1,00). Khi vượt số core và $\mathrm{OW} \ge 7$, MCS kém công bằng nhất trong mười hai lock
(0,21--0,30 ở 16--20 thread, so với 0,78--0,81 của lock kế tiếp). Nhóm cho rằng một thread MCS bị lấy core giữa
\texttt{release} và \texttt{acquire} không nằm trong hàng đợi, nên các thread còn chạy lấy lock nhiều
lượt hơn. Khi $\mathrm{OW}$ lớn, khoảng giữa \texttt{release} và \texttt{acquire} dài, nên điều này xảy
ra thường hơn.
% nguồn: lock.csv (mcs N=1009, độ công bằng thr_min_s/thr_max_s): p<=10 mọi OW 0,84..1,00; p=12/16/20: OW=0 1,00/0,99/0,99,
% OW=1 1,00/0,96/0,95, OW=3 1,00/0,84/0,94; OW=7 p=16/20 0,30/0,21; OW=15 p=16/20 0,29/0,26 (rlock_analyze.py cùng số).
% Lock kế tiếp ở OW=7/15, p=16/20: 0,78..0,81.

Khi vượt số core và $\mathrm{OW} \ge 3$, \PCname{ttas_bo_sleep} chậm hơn TTAS 6,1--6,8 lần ở
$\mathrm{OW} = 3$ với 16--20 thread, 2,0--4,3 lần ở $\mathrm{OW} = 7$ và 1,6--1,8 lần ở
$\mathrm{OW} = 15$ (12--20 thread). Nhóm giả thuyết rằng ở biến thể này một phần thread chờ quay đọc và
chiếm core, phần khác ngủ. Thread vừa thức dậy sau khi ngủ có thể giành core của chủ lock. Trong lúc chủ
lock không chạy, các thread quay đọc dùng hết lượt của mình. \PCname{ttas_bo_yield}, chỉ khác ở chỗ
nhường core thay vì ngủ, không chậm hơn TTAS (1,97~s so với 2,42~s ở 20 thread, $\mathrm{OW} = 3$, thời
gian quy về cùng số vòng lặp). \PCname{tas_bo_sleep}, nơi mọi thread chờ đều ngủ, cũng không bị. Hai
số đo này phù hợp với giả thuyết.
% nguồn: scripts/report_numbers/rlock_analyze.py tỉ số ttas_bo_sleep/ttas N=1009: OW=3 p=16/20 6,13/6,78;
% OW=7 (độ chậm 16,6/3,89, 16,1/4,53, 11,5/5,85) p=12/16/20 4,27/3,55/1,97; OW=15 (5,63/3,28, 9,03/4,93, 9,1/5,82) 1,72/1,83/1,56.
% tables.tex locktime-ow3 p=20: ttas_bo_yield 1,97, ttas 2,42, ttas_bo_sleep 16,4; tas_bo_sleep
% 4,27 so với tas 6,20; p=16 ttas_bo_sleep 10,4, ttas 1,70.

\emph{Độ vững.} Kết quả của \PCname{ttas_bo_sleep} khi vượt số core phân tán mạnh giữa các lần lặp. Ở
16 thread và $\mathrm{OW} = 3$, thời gian quy về cùng số vòng lặp đi từ 1,0 tới 18,0~s, với trung vị
10,4~s. Đổi lệnh chờ \texttt{isb} thành \texttt{yield} không tạo
khác biệt vượt độ phân tán.
% nguồn: tables.tex locktime-ow3 p=16/20: ttas_bo_sleep 10,4/16,4, ttas 1,70/2,42; lock.csv ttas_bo_sleep
% p=16 OW=3: iters 21\,479, thô 0,259..4,510 s, quy về I = 85\,916 (x4,0): 1,04..18,04 s, trung vị 2,609 x 4 = 10,4 s.
% lock_yield.csv vs lock.csv: trung vị tỉ số 0,86..1,09, min..max 0,69..2,17.
